\documentclass[11pt]{article}

\usepackage[T1]{fontenc}
\usepackage[utf8]{inputenc}
\usepackage{lmodern}
\usepackage[margin=1in]{geometry}
\usepackage{microtype}
\usepackage{amsmath,amssymb,mathtools}
\usepackage{booktabs,tabularx,array}
\usepackage{enumitem}
\usepackage{xcolor}
\usepackage{hyperref}

\hypersetup{
  colorlinks=true,
  linkcolor=blue!55!black,
  citecolor=blue!55!black,
  urlcolor=blue!55!black
}

\newcommand{\Spec}{\operatorname{Spec}}
\newcommand{\Tr}{\operatorname{Tr}}
\newcommand{\cA}{\mathcal{A}}
\newcommand{\cC}{\mathcal{C}}
\newcommand{\cP}{\mathcal{P}}
\newcommand{\cS}{\mathcal{S}}
\newcommand{\R}{\mathbb{R}}

\title{Generative Relational Coherence\\
\large Toward a Higher-Level Ontology of Physics}
\author{}
\date{}

\begin{document}

\maketitle

\begin{abstract}
The representational fragmentation exposed by the Yang--Mills mass-gap problem
points beyond a philosophy of mathematical description toward a possible
ontology of physics.  The appropriate generalization is not that the universe
is simply information.  Information presupposes distinguishable alternatives,
physical relations, constraints, probabilities, and rules governing how
distinctions propagate.  A more fundamental proposal is \emph{generative
relational coherence}: the universe is a dynamically self-consistent structure
of relations and possible processes in which local constraints generate stable
global invariants across scale, context, and representation.  Fields,
particles, regulators, renormalization, continuum limits, and mass gaps can then
be understood as different kinds of presentation, transformation, coherence,
and rigidity within this structure.  This essay develops that proposal and
offers a precise answer to the question of what a universe is from this point
of view.
\end{abstract}

\section{From the Yang--Mills problem to a general ontology}

The conceptual statement arising from the Yang--Mills mass-gap problem is:
\begin{quote}
The Yang--Mills mass-gap problem is fundamentally the problem of generating and
transporting infrared spectral rigidity across a chain of representations whose
natural variables separately expose locality, gauge invariance, positivity,
scale dependence, and spectral structure.
\end{quote}

This formulation can be generalized:
\begin{quote}
Fundamental physical problems concern how local relational constraints generate
global forms of rigidity, and how those rigidities remain identifiable across
representations that separately expose different aspects of physical reality.
\end{quote}

Schematically,
\begin{equation}
\boxed{
\text{local relational rules}
\longrightarrow
\text{collective organization}
\longrightarrow
\text{global invariant rigidity}.}
\label{eq:general-pattern}
\end{equation}
Each arrow can involve the composition of many degrees of freedom, changes of
scale, elimination of redundant variables, singular limits, emergence of new
effective objects, and translation between mathematical representations.

For Yang--Mills theory the pattern becomes
\begin{equation}
  \text{local non-Abelian interaction}
  \longrightarrow
  \text{organized quantum vacuum}
  \longrightarrow
  \text{infrared spectral rigidity}.
  \label{eq:ym-pattern}
\end{equation}
Related patterns appear when microscopic reversibility generates macroscopic
irreversibility, local quantum correlations generate topological order,
microscopic interactions generate thermodynamic phases, local causal processes
generate effective spacetime geometry, or quantum many-body dynamics generates
quasiparticles.

The higher proposition is therefore:
\begin{quote}
Physics studies the generation, transformation, and preservation of relational
rigidity.
\end{quote}

\section{What is more fundamental than information?}

The statement that the universe is information is suggestive but incomplete.
Information already presupposes a space of alternatives, criteria of
distinguishability, physical observables, probability assignments, and rules
relating one subsystem to another.  A more fundamental hierarchy is
\begin{equation}
\boxed{
\text{relations}
\longrightarrow
\text{constraints}
\longrightarrow
\text{distinctions}
\longrightarrow
\text{information}
\longrightarrow
\text{stable objects}.}
\label{eq:ontological-hierarchy}
\end{equation}

\subsection{Relations}

Relations express how physical possibilities depend on one another: causal
influence, spatial adjacency, interaction, entanglement, correlation,
composition, and exclusion.

\subsection{Constraints}

Constraints determine which joint relational configurations are physically
possible.  Dynamical equations, conservation laws, causal order, positivity,
symmetry, gauge constraints, and composition rules belong to this level.

\subsection{Distinctions}

A distinction is a physically realizable contrast between possibilities.  It
does not require a conscious observer.  A detector, field, environment, or
second subsystem may distinguish two physical states through interaction.

\subsection{Information}

Information quantifies distinguishability within a physical context.  If
$\omega$ and $\nu$ are states available in a context $c$, one may write
schematically
\begin{equation}
  D_c(\omega,\nu)
  =
  \sup_{O\in\cA(c)}
  \left|\omega(O)-\nu(O)\right|,
  \label{eq:distinguishability}
\end{equation}
where $\cA(c)$ is the class of observables accessible in that context.  The
available information depends on the physically realizable distinctions, not
merely on an abstract list of hidden values.

\subsection{Stable objects}

An object is a pattern of relational structure that remains identifiable under
an appropriate class of transformations.  Information is therefore not an
independent substance from which relations are built.  It is a quantitative
measure defined on a prior structure of constrained relational possibilities.

\section{Objects as invariants of transformation}

Within this ontology, an object is not fundamentally a self-contained bearer of
intrinsic properties.  It is a stable structure persisting across transformations
of description, scale, localization, and interaction.

A particle, for example, can appear as
\begin{itemize}[leftmargin=2em]
  \item a pole or characteristic singularity of a correlation function;
  \item an irreducible representation of spacetime symmetries;
  \item a spectral sector of a Hamiltonian;
  \item a localized detector response;
  \item an asymptotic scattering state;
  \item an excitation of an underlying field.
\end{itemize}
None of these manifestations alone is simply the particle.  Where the relevant
equivalence and reconstruction theorems apply, the particle is the coherent
invariant relating them.

This suggests the definition
\begin{equation}
  \boxed{
  \text{physical object}
  =
  \text{stable equivalence class of relational manifestations}.}
  \label{eq:object-definition}
\end{equation}

Different entities exhibit different forms of stability.

\begin{table}[ht]
\centering
\small
\renewcommand{\arraystretch}{1.2}
\begin{tabularx}{0.92\textwidth}{>{\raggedright\arraybackslash}p{0.30\textwidth}X}
\toprule
\textbf{Entity} & \textbf{Characteristic form of stability} \\
\midrule
Particle & Spectral and symmetry stability \\
Field & Local response and compositional stability \\
Phase of matter & Stability under perturbation and coarse-graining \\
Topological excitation & Stability under local deformation \\
Spacetime event & Stability within causal and localization relations \\
Conservation law & Invariance under admissible dynamics \\
Mass gap & Stability of the vacuum against arbitrarily low-energy excitation \\
\bottomrule
\end{tabularx}
\caption{Physical entities as distinct forms of relational stability.}
\label{tab:stabilities}
\end{table}

Objects can consequently be real without possessing one privileged
mathematical presentation.

\section{Reinterpreting the Yang--Mills statements}

The earlier propositions can now be given an ontological reading.

\subsection{Gauge fields are coordinates on physical information}

A gauge field $A_\mu$ is not directly physical because gauge-related
configurations represent the same physical situation.  Its role resembles a
coordinate system on an underlying relational structure.  The physical content
is carried by such structures as curvature, holonomies, Wilson loops,
gauge-invariant correlations, responses to admissible probes, and the physical
state space.

Thus, schematically,
\begin{equation}
  A_\mu
  \quad\text{is a presentation, while}\quad
  [A_\mu]_{\mathrm{gauge}}
  \quad\text{carries physical identity}.
  \label{eq:gauge-identity}
\end{equation}
Even the gauge-equivalence class need not be the final ontology.  It may itself
be a representation of a more directly observable-relational structure.

\subsection{Regulators are temporary realizations}

A lattice is more than a computational approximation.  It is a concrete
realization of a family of relational constraints.  The continuum theory is
physical only if a coherent structure survives as
\begin{equation}
  a\longrightarrow0,
  \qquad
  L\longrightarrow\infty.
\end{equation}
The lattice variables may disappear while their invariant observable relations
converge.  The regulator is therefore a temporary carrier of physical
structure.

\subsection{Particles are spectral interpretations}

A particle need not be a microscopic primitive.  It can be a stable spectral
organization of the theory:
\begin{equation}
  \Spec(H)
  =
  \{0\}
  \cup
  \{\text{isolated massive sectors}\}
  \cup
  \{\text{spectral continua}\}.
\end{equation}
In a strongly interacting theory the natural physical excitations may be very
different from the elementary variables in the classical action.  In pure
Yang--Mills theory, gauge potentials appear in the local equations, whereas
gauge-invariant composite excitations are expected in the physical spectrum.

\subsection{Renormalization transforms descriptive scale}

Renormalization does not ordinarily represent physical time evolution.  It
changes which distinctions are retained at a selected resolution.  A
coarse-graining transformation
\begin{equation}
  \mathcal R_{\ell\to\ell'}:
  \mathfrak R_\ell
  \longrightarrow
  \mathfrak R_{\ell'}
  \label{eq:coarse-graining}
\end{equation}
forgets some microscopic distinctions while preserving selected long-distance
predictions.  Universal quantities are those that remain stable despite this
loss of microscopic descriptive detail.

\subsection{The continuum is a coherence condition}

The continuum need not be understood merely as an infinitely fine lattice.  It
is the coherent limiting object represented by a compatible family of
regulated descriptions.  If $\omega_a$ is the state at cutoff $a$, then for
appropriate observables one requires
\begin{equation}
  \lim_{a\to0}\omega_a(O_a)=\omega(O),
  \label{eq:continuum-coherence}
\end{equation}
together with compatibility across observable products, symmetries,
localization, positivity, and scale transformation.  The continuum is thus a
consistency structure across a family of approximations rather than another
microscopic picture.

\subsection{The mass gap is infrared rigidity}

The mass gap states that the vacuum cannot be changed into a distinct physical
excitation at arbitrarily small positive energetic cost:
\begin{equation}
  H\geq\Delta(1-P_0),
  \qquad \Delta>0,
  \label{eq:mass-gap}
\end{equation}
where $P_0$ projects onto the vacuum sector.  Ontologically, this says that the
vacuum is rigid against sufficiently long-wavelength physical deformation.

The corresponding decay
\begin{equation}
  \langle O(x)O(0)\rangle_c
  \sim e^{-\Delta|x|}
  \label{eq:clustering}
\end{equation}
means that local physical distinctions do not propagate as indefinitely
persistent low-energy correlations.  The vacuum has a finite relational memory
length.

\section{The universe as a coherent family of contexts}

The ontology can be made more precise by introducing physical contexts.  Let
\begin{equation}
  c=(\Omega,\ell,B,P,R),
  \label{eq:context}
\end{equation}
where $\Omega$ is a spacetime region, $\ell$ a scale or resolution, $B$ a set of
boundary conditions, $P$ a class of physical probes, and $R$ a mathematical
presentation.

To every context $c$, assign
\begin{equation}
  \mathfrak U(c)
  =
  \bigl(
    \cA(c),
    \cS(c),
    \cP(c),
    \cC(c)
  \bigr),
  \label{eq:context-assignment}
\end{equation}
where
\begin{itemize}[leftmargin=2em]
  \item $\cA(c)$ is the observable algebra;
  \item $\cS(c)$ is the state space;
  \item $\cP(c)$ is the space of possible physical processes;
  \item $\cC(c)$ records causal, compositional, and compatibility constraints.
\end{itemize}

Contexts are related by transformations
\begin{equation}
  \Phi_{c\to d}:\mathfrak U(c)\longrightarrow\mathfrak U(d).
  \label{eq:context-map}
\end{equation}
These transformations may express localization, inclusion of regions, change
of resolution, coarse-graining, gauge quotient, change of coordinates, duality,
continuum reconstruction, or physical time evolution.

A universe is not any single description $\mathfrak U(c)$.  It is the
structured totality
\begin{equation}
  \mathbb U
  =
  \left(
    \{\mathfrak U(c)\}_{c\in\mathsf C},
    \{\Phi_{c\to d}\},
    \{\eta_{c,d,e}\}
  \right),
  \label{eq:universe-totality}
\end{equation}
where $\eta_{c,d,e}$ expresses coherence between alternative transformation
paths.  For example, transforming from $c$ to $e$ directly should agree,
exactly or within a controlled approximation, with transforming through $d$:
\begin{equation}
  \Phi_{c\to e}
  \simeq
  \Phi_{d\to e}\circ\Phi_{c\to d}.
  \label{eq:path-coherence}
\end{equation}
The universe is not one member of this family.  It is the coherent organization
of the entire family.

\section{Global coherence without predetermined values}

This ontology does not say that every observable has a predetermined value or
that the universe is a hidden classical database.  Quantum theory obstructs
such a naive global value assignment.  The compatible global object is instead
a state or probability structure assigning expectations to physically
admissible observables:
\begin{equation}
  \omega_c:\cA(c)\longrightarrow\mathbb C.
\end{equation}

Compatibility requires, schematically,
\begin{equation}
  \omega_d(O_d)
  =
  \omega_c\!\left(\Phi_{c\to d}^{*}O_d\right)
  \label{eq:state-compatibility}
\end{equation}
whenever $c$ and $d$ are exact descriptions of the same physical content.  The
universe may therefore be globally coherent without being a catalogue of
simultaneously definite local facts.

\section{Objectivity as certified coherence}

The contextual character of representation does not entail relativism.
Different representations are not automatically equally adequate.  A
representation can fail because it loses positivity, changes the physical
observable algebra, fails to reconstruct causal locality, eliminates
nontrivial interactions, has the wrong continuum limit, destroys spectral
structure, or introduces an uncontrolled approximation.

Objectivity is therefore not independence from every representation.  It is
invariance under the admissible transformations connecting valid
representations:
\begin{equation}
\boxed{
  \text{objective physical content}
  =
  \text{content stable under certified changes of representation}.}
  \label{eq:objectivity}
\end{equation}
This is a form of structural realism with a generative component.  Reality
contains both invariant structures and lawful processes through which new
stable structures emerge.

\section{Why transformations must remain typed}

Not every transformation represents physical change.  The ontology must keep
distinct transformations distinct.

\begin{table}[ht]
\centering
\small
\renewcommand{\arraystretch}{1.2}
\begin{tabularx}{\textwidth}{>{\raggedright\arraybackslash}p{0.28\textwidth}X}
\toprule
\textbf{Transformation} & \textbf{Physical or representational meaning} \\
\midrule
Gauge transformation & Same physical state, different presentation \\
Coordinate transformation & Same structure, different coordinates \\
Exact duality & Same physical content, different ontology of variables \\
Coarse-graining & Reduced resolution and controlled information loss \\
Restriction to a region & Partial physical perspective \\
Continuum limit & Construction of a limiting theory \\
Time evolution & Actual physical change of state \\
Measurement interaction & Physical correlation and outcome formation \\
\bottomrule
\end{tabularx}
\caption{Different transformation types must not be conflated.}
\label{tab:transformations}
\end{table}

Representational dynamics is a meta-level theory of how descriptions
transform.  Physical dynamics is one distinguished class of transformations
within it.  A precise formalization may therefore require a double or higher
categorical structure: vertical arrows for equivalent descriptions, horizontal
arrows for restriction, coarse-graining, and physical processes, and higher
cells certifying that different transformation paths are coherent.

\section{Laws as coherence constraints}

In an object-centered ontology, laws primarily describe how things move or
change.  In a relational-coherence ontology, laws have a broader function.
They determine
\begin{itemize}[leftmargin=2em]
  \item which local relational configurations are admissible;
  \item how systems and processes compose;
  \item which transformations count as physically equivalent;
  \item how descriptions at different scales correspond;
  \item which quantities are conserved;
  \item which emergent structures remain stable;
  \item how local information constrains global possibilities.
\end{itemize}

A physical law is therefore not merely a differential equation.  It is a
coherence condition across a family of possible physical contexts.  The
Yang--Mills action remains essential, but it is one compact local generator
within a larger structure including gauge equivalence, observables, states,
renormalization, positivity, limiting constructions, and spectral consequences.

\section{What is a universe from this point of view?}

A concise answer is:
\begin{quote}
A universe is a dynamically self-contained and globally coherent system of
relational possibilities, physical constraints, and composable processes,
whose valid descriptions form a compatible family across localization, scale,
and representation, and whose objects are the stable invariants generated
within that family.
\end{quote}

Schematically,
\begin{equation}
\boxed{
  \text{Universe}
  =
  \text{relations}
  +
  \text{constraints}
  +
  \text{processes}
  +
  \text{coherence}
  +
  \text{emergent rigidity}.}
  \label{eq:universe-definition}
\end{equation}

From this perspective,
\begin{itemize}[leftmargin=2em]
  \item spacetime is the organization of localization and causal
        composability;
  \item matter is a family of stable excitation structures;
  \item particles are persistent spectral patterns;
  \item fields encode possible local responses and correlations;
  \item symmetries classify equivalence and conservation;
  \item information measures physically available distinctions;
  \item laws regulate composition and compatibility;
  \item renormalization relates descriptions at different resolutions;
  \item objects are structures stable under relevant transformations;
  \item emergence is the generation of new invariants under composition and
        scale change;
  \item the universe is the maximal coherent physical process containing these
        relations.
\end{itemize}

\section{The strongest generalized proposition}

The original Yang--Mills statement can finally be generalized as follows:
\begin{quote}
A physical universe is a generative relational structure in which local
constraints produce global invariant rigidities, while lawful transformations
connect the partial representations through which those rigidities become
observable, calculable, and provable.
\end{quote}

The Yang--Mills mass gap is a particularly pure test case.  It asks whether a
local, gauge-redundant, scale-free presentation generates a gauge-invariant
infrared rigidity that survives every required transformation into the
continuum physical theory.

This is why the problem is philosophically central.  It tests whether physics
can explain not merely the behavior of already constituted objects, but the
emergence of stable physical objecthood itself from relational quantum
structure.

\end{document}
